% ===================================================================
\section{Resultados}

\diapo{18}{Validación numérica}{45 s}
\begin{decir}
Antes de interpretar resultados hay que confiar en el código. Cada componente se
comparó con un resultado conocido.

La integración coincide con la solución exacta del sistema lineal con error menor que
$10^{-6}$. La redistribución conserva el déficit total con precisión de $10^{-10}$. La
sincronización reproduce el valor teórico de Kuramoto. El modelo multicapa con una sola
capa coincide exactamente con la versión anterior. Y la coherencia estimada desde
transacciones se correlaciona por encima de 0.8 con la real.

El conjunto de datos se regenera idéntico bit a bit. En total el proyecto tiene 127
pruebas automáticas, y todas pasan.
\end{decir}
\begin{terminos}
\begin{description}
\item[RK4.] Método de Runge–Kutta de cuarto orden: evalúa la derivada cuatro veces por
paso y combina los resultados. Error global proporcional a $\Delta t^4$. Integrador de
las tareas 1–5.
\item[$e^{Mt}x_0$.] Solución exacta de $\dot x=Mx$ con condición inicial $x_0$; $e^{Mt}$
es la exponencial de una matriz.
\item[Varianza de Ornstein–Uhlenbeck.] Con solo $-\alpha x+\sigma\xi$, la varianza
estacionaria es $\sigma^2/(2\alpha)$. La simulación la reproduce dentro del 10~\% y
valida el ruido.
\item[Bit a bit.] Dos ejecuciones producen archivos idénticos byte por byte.
\item[Prueba automática.] Función que ejecuta una parte del código y comprueba su
resultado; se corre con \texttt{pytest}.
\end{description}
\end{terminos}

\diapo{19}{Carteras simuladas}{45 s}
\begin{decir}
Estas son las carteras del experimento. En el régimen nulo hay unos 40\,000 clientes
modelados, con una tasa de impago a doce meses de 5.9~\%. En el régimen con red hay unos
37\,000, porque más clientes caen durante la observación, y la tasa de impago sube a
9.6~\%. La diferencia es el efecto de la red. El índice de amplificación está entre 0.09
y 0.17: la red amortigua, no está cerca de la inestabilidad.

La figura muestra la coherencia estimada por grupo. Los empleados formales están
sincronizados con su nómina. Los informales reciben ingresos de pagadores con ritmos
distintos y su coherencia está cerca de cero. En el régimen con red, los impagos por
umbral, 2\,335, superan a los idiosincráticos, 1\,264.
\end{decir}
\begin{terminos}
\begin{description}
\item[Clientes modelados.] Clientes vivos al final de la observación, sumados en las 20
carteras.
\item[Diagrama de caja.] La caja va del primer al tercer cuartil; la línea interior es la
mediana; los bigotes muestran la dispersión.
\item[Impago por umbral.] Causado por la dinámica: el déficit superó el colchón.
\item[Impago idiosincrático.] Programado por la PD propia, independiente de la red.
\end{description}
\end{terminos}

\diapo{20}{Desempeño predictivo por nivel}{1 min}
\begin{decir}
Este es el resultado principal. En el régimen con red, con regresión logística, el AUC
pasa de 0.769 con el score clásico a 0.784 al agregar la red estática, y a 0.795 al
agregar la dinámica. Con XGBoost el patrón es el mismo y llega también a 0.795. Las
demás métricas mejoran en la misma dirección: la PR-AUC, que es la métrica relevante
cuando los impagos son pocos, sube de 0.277 a 0.323.

En el régimen nulo, en la parte izquierda de la figura, el AUC es más alto, 0.81, pero
los tres niveles son iguales. Es más alto porque allí casi todos los impagos dependen
solo de la probabilidad propia de cada cliente, que es justamente lo que mide el score
clásico. Los tres niveles son iguales porque la red no tiene nada que aportar.
\end{decir}
\begin{terminos}
\begin{description}
\item[Niveles 1, 2 y 3.] Clásico, clásico más red estática, y clásico más red estática
más red dinámica.
\item[PR-AUC.] Promedio de la precisión sobre todos los niveles de exhaustividad. Un
modelo al azar obtiene la tasa de impago, aquí cerca de 0.096.
\item[Precisión.] De los clientes marcados como riesgosos, fracción que cae.
\item[Exhaustividad (\emph{recall}).] De los clientes que caen, fracción marcada como
riesgosa.
\item[KS.] Estadístico de Kolmogorov–Smirnov: máxima distancia entre las distribuciones
acumuladas del score de quienes caen y de quienes no. Es una métrica habitual en banca.
\item[Brier.] Error cuadrático medio de la probabilidad, $\frac1n\sum(p_i-y_i)^2$.
Menor es mejor.
\item[ECE.] Error de calibración esperado: diferencia media entre la probabilidad
predicha y la frecuencia observada, por tramos de probabilidad.
\item[Fuera de muestra.] Métricas calculadas sobre carteras que el modelo no vio.
\end{description}
\end{terminos}
\begin{preguntas}
\pr{¿Por qué el AUC del régimen nulo es mayor?}{En el nulo el impago depende casi solo
de la PD de cada cliente, que es función de las variables clásicas. En el régimen con
red, parte del impago viene de la red y el score clásico no lo ve. Por eso el score
clásico pierde AUC con red y las variables de red recuperan una parte.}
\end{preguntas}

\diapo{21}{Mejoras pareadas y control nulo}{1 min}
\begin{decir}
Para saber si las diferencias son reales, comparamos los niveles dentro de cada
partición. La figura muestra cada mejora con su intervalo de confianza del 95~\%; en
verde, las que están por encima de cero.

En el régimen con red, la dinámica mejora a la red estática en 0.011 de AUC con
regresión logística, con intervalo entre 0.006 y 0.016, y en 0.014 con XGBoost. La
mejora es positiva en las cinco particiones. El intervalo no incluye el cero.

En el régimen nulo, la mejora está entre $-0.002$ y $+0.001$. Este es el control: cuando
el mecanismo no existe, el método no encuentra nada. Eso descarta que la mejora venga de
una fuga de información o de que el modelo tenga más variables.
\end{decir}
\begin{terminos}
\begin{description}
\item[$\Delta\mathrm{AUC}$.] AUC del nivel superior menos AUC del nivel de comparación,
en la misma partición.
\item[Gráfico de bosque.] Cada fila es una comparación: el punto es la media y la barra
el intervalo de confianza.
\end{description}
\end{terminos}
\begin{preguntas}
\pr{¿Una mejora de 0.01 en AUC es relevante?}{En riesgo de crédito sí. El Gini, que usan
los bancos, es $2\,\mathrm{AUC}-1$, así que la mejora es de 0.022 puntos de Gini sobre un
modelo que ya incluye la red estática. Sobre el score clásico, la mejora total con
XGBoost es de 0.031 de AUC. Además, el valor depende de parámetros ilustrativos: lo que
el experimento establece es que la mejora existe, es estable y se anula sin el
mecanismo.}
\pr{¿Más variables no siempre mejoran el AUC?}{No fuera de muestra: variables sin
información agregan ruido. El régimen nulo lo muestra: con las mismas variables, la
mejora es cero.}
\end{preguntas}

\diapo{22}{Respuesta a la intensidad del mecanismo}{1 min}
\begin{decir}
Este experimento adicional responde de dónde viene la mejora. Variamos $\gamma$, el
parámetro que convierte el desfase en déficit, y dejamos todo lo demás igual.

Con $\gamma$ igual a cero, el contagio y la pérdida de ingreso siguen activos, y aun así la
dinámica no aporta nada. La mejora aparece solo cuando el desfase genera déficit, y crece
de manera monótona: 0.001, 0.008 y 0.029 de AUC con $\gamma$ igual a 0.25, 0.5 y 1.

Conclusión: la información que aportan las variables dinámicas proviene de la
desincronización, no del contagio. Es una relación dosis-respuesta: más mecanismo, más
señal.
\end{decir}
\begin{terminos}
\begin{description}
\item[Sensibilidad.] Variación del resultado ante cambios en un parámetro.
\item[Diseño del experimento.] Regresión logística, 10 carteras por valor de $\gamma$,
misma validación; barras de error con IC 95~\%.
\item[Monótona.] Que siempre crece (o siempre decrece) al aumentar el parámetro.
\item[Dosis-respuesta.] Término de epidemiología: el efecto crece con la intensidad de la
causa. Es un criterio clásico para atribuir un efecto a una causa.
\end{description}
\end{terminos}

\diapo{23}{Interpretación y relación sincronía--impago}{1 min}
\begin{decir}
A la izquierda, la importancia de cada variable en el modelo híbrido. Las clásicas,
en gris, siguen siendo las más importantes: endeudamiento, ingreso y moras. Las de red
estática, en azul, y las dinámicas, en rojo, aportan información adicional. Entre las
dinámicas destacan la variabilidad de la coherencia, su nivel medio y el desfase. El
modelo híbrido complementa al score; no lo reemplaza.

A la derecha, la tasa de impago según la coherencia estimada. En el régimen con red, el
impago pasa de 4.4~\% en los clientes bien sincronizados a entre 11 y 14~\% en los
desincronizados. En el régimen nulo también hay una diferencia, más pequeña, porque los
informales son más riesgosos y además están menos sincronizados. Por eso esta gráfica
por sí sola no prueba nada: la prueba es la comparación de modelos de las diapositivas
anteriores.
\end{decir}
\begin{terminos}
\begin{description}
\item[SHAP.] \emph{SHapley Additive exPlanations} (Lundberg y Lee, 2017). Reparte la
predicción de un cliente entre sus variables usando los valores de Shapley de la teoría
de juegos: la contribución de cada variable promediada sobre todos los órdenes en que se
pueden agregar las variables. La suma de las contribuciones más un valor base da la
predicción.
\item[Importancia media $|\mathrm{SHAP}|$.] Promedio sobre los clientes del valor
absoluto de la contribución de cada variable, en unidades de log-odds.
\item[Log-odds.] $\ln\frac{p}{1-p}$, la escala en que el modelo suma contribuciones.
\item[Variable de confusión.] Variable que influye a la vez en la causa aparente y en el
efecto. Aquí, la informalidad.
\end{description}
\end{terminos}
\begin{preguntas}
\pr{¿El modelo está calibrado?}{Sí. En la curva de calibración ambos modelos siguen la
diagonal y el híbrido separa mejor el decil de mayor riesgo: 35~\% de impago observado
frente a 30~\% del clásico. La figura está en la guía completa y en la presentación v1.}
\end{preguntas}

\diapo{24}{Escenario: quiebra de un empleador}{1 min}
\begin{decir}
El modelo también sirve para pruebas de estrés. Simulamos la quiebra del mayor empleador
de la cartera, con 74 empleados, en el primer mes, y promediamos cinco semillas.

En el régimen con red, el 12~\% de los empleados cae de inmediato, porque pierde todo su
salario de un momento a otro. A los seis meses el impago acumulado de esos empleados es
13~\%, frente a 7~\% sin quiebra. En el régimen nulo la quiebra no tiene efecto, como
corresponde.

La aplicación directa para un banco es medir cuánto riesgo concentra en los empleados de
una misma empresa.
\end{decir}
\begin{terminos}
\begin{description}
\item[Prueba de estrés.] Simulación de un escenario adverso extremo para medir pérdidas.
\item[Choque de quiebra.] Impago forzado de la firma en $t_0=1$ mes: deja de pagar
nómina, proveedores y cuota.
\item[Semilla.] Número que inicializa el generador aleatorio; misma semilla, mismo
resultado.
\end{description}
\end{terminos}

\diapo{25}{Alcance e interpretación}{45 s}
\begin{decir}
Para cerrar los resultados, separo lo que el diseño respalda de lo que no.

Respalda tres afirmaciones. Si el mecanismo existe, es detectable con las transacciones
de que ya dispone un banco. La dinámica aporta información que la red estática no tiene.
Y el aporte se anula sin el mecanismo y crece con su intensidad.

No respalda que el mecanismo exista en una cartera real: los datos son sintéticos y los
parámetros ilustrativos. Además, supusimos que el banco ve toda la red de sus clientes.
\end{decir}
\begin{preguntas}
\pr{¿Qué haría falta para usarlo en un banco?}{Aplicar el mismo pipeline a una cartera
real anonimizada, bajo la Ley 1581 de 2012 de protección de datos personales. El código
ya recibe un registro de transacciones con la misma estructura.}
\end{preguntas}

% ===================================================================
\section{Implementación}

\diapo{26}{Arquitectura y reproducibilidad}{30 s}
\begin{decir}
El software está organizado en cuatro capas. La interfaz es una aplicación web. La capa
de servicio prepara las tablas que la interfaz muestra. El núcleo contiene el modelo y los
algoritmos. La configuración y los datos están separados del código. Cada capa solo
depende de la de abajo.

Todos los parámetros están en archivos de configuración, las semillas son explícitas, las
tablas tienen esquemas validados, hay 127 pruebas, integración continua y una imagen de
Docker que ejecuta la aplicación igual en cualquier equipo.
\end{decir}
\begin{terminos}
\begin{description}
\item[Streamlit.] Biblioteca de Python que convierte un programa en una aplicación web
interactiva.
\item[YAML.] Formato de texto para configuración, legible por personas.
\item[Esquema validado.] Lista de columnas, tipos y rangos que toda tabla debe cumplir;
si no se cumple, el programa se detiene.
\item[Integración continua.] Ejecución automática de pruebas en cada cambio del código
(GitHub Actions).
\item[Docker.] Herramienta que empaqueta el programa con todas sus dependencias en una
imagen que se ejecuta igual en cualquier computador.
\end{description}
\end{terminos}

\diapo{27}{Flujo de datos de la tarea 6}{30 s}
\begin{decir}
Este es el recorrido completo de la tarea 6: generar la cartera, simular el modelo,
producir las transacciones, estimar la red y las fases, construir las variables, validar
y comparar los modelos. Un solo comando ejecuta todo el flujo en unos 20 segundos.
\end{decir}
\begin{terminos}
\begin{description}
\item[\texttt{uv}.] Gestor de entornos y dependencias de Python.
\item[\texttt{pulso credit}.] Comando del proyecto que ejecuta el experimento de la
tarea 6 y guarda resultados y modelos finales.
\end{description}
\end{terminos}

\diapo{28}{Aplicación web}{30 s}
\begin{decir}
Todo lo presentado se puede explorar en la aplicación. A la izquierda, el simulador de
choques de las tareas 1 a 4: se elige una entidad, se aplica un choque y se ve la
cascada. A la derecha, la comparación de los tres niveles de la tarea 6 y la
probabilidad de impago por cliente, con su explicación. El manual de la aplicación
describe cada página.
\end{decir}

% ===================================================================
\section{Conclusiones}

\diapo{29}{Conclusiones}{1 min}
\begin{decir}
Primero, un único modelo dinámico integra contagio, redistribución, sincronización y
cascadas, y cada componente se valida contra resultados analíticos.

Segundo, el índice $R$ y el espectro del sistema dicen cuándo la red amortigua y cuándo
amplifica un choque.

Tercero, el aprendizaje automático predice el colapso del sistema con AUC de 0.98 en
milisegundos.

Cuarto, y es el resultado central: en una cartera multicapa, la sincronización de los
ciclos de pago, estimada solo con transacciones, mejora la predicción del impago sobre la
red estática. No tiene efecto en el control nulo y crece con la intensidad del
mecanismo.

Quinto, el software es reproducible, está probado y funciona como aplicación web.
\end{decir}

\diapo{30}{Trabajo futuro}{30 s}
\begin{decir}
El paso más importante es la validación con datos reales de un banco, anonimizados.
También queda pendiente una revisión sistemática de la literatura. En el modelo, faltan
las fechas de cuota como ritmo impuesto desde fuera, la retroalimentación del déficit
sobre el ritmo de pago, y el caso en que un banco solo ve a sus clientes, frente al
operador de Bre-B, que ve la red completa.
\end{decir}
\begin{terminos}
\begin{description}
\item[Revisión sistemática.] Búsqueda bibliográfica con protocolo explícito en bases como
Scopus y Web of Science.
\item[Forzamiento externo.] Ritmo impuesto desde fuera del sistema, por ejemplo la fecha
fija de la cuota.
\item[Observabilidad parcial.] Situación en la que solo se ve una parte de la red.
\end{description}
\end{terminos}

\diapo{31}{Referencias principales}{15 s}
\begin{decir}
Estas son las referencias principales. La bibliografía completa está en el repositorio.
Muchas gracias; quedo atento a sus preguntas.
\end{decir}
